\chapter{Porting to a New Platform}\label{ch:porting}

A port is one source file, sometimes two, that implements the functions in
\file{include/eip_hal.h}. The core calls these functions directly. There is
no function table. This chapter tells you what each function must do and how
to check a new port.

\section{Select a Starting Point}

Start from the port that is nearest to your platform.
Table~\ref{tab:portstart} shows which port to use.

\begin{table}[htbp]
\centering
\caption{Starting points for a new port}\label{tab:portstart}
\small
\begin{tabularx}{\linewidth}{@{}P{0.32\linewidth}P{0.30\linewidth}L@{}}
\toprule
\textbf{Your platform has} & \textbf{Start from} & \textbf{You write} \\
\midrule
BSD sockets (\code{socket}, \code{select}, \code{fcntl}): lwIP with \code{LWIP_SOCKET}, Zephyr, Linux, FreeRTOS+TCP with its BSD layer & \file{ports/bsd_sockets/eip_hal_bsd_sockets.c}, without changes & A platform file like \file{ports/esp-idf/eip_hal_esp_idf.c}: clock, lock, interfaces, and task \\
The lwIP raw API, \code{NO_SYS=1}, no socket layer: STM32Cube bare metal, SAM, other boards like ClearCore & \file{ports/lwip_raw/eip_hal_lwip_raw.c}, without changes & A platform file like \file{ports/clearcore/eip_hal_clearcore.cpp}: clock, yield (your stack poll), and interfaces \\
A vendor socket API that is neither (NetBurner) & \file{ports/netburner/eip_hal_netburner.cpp} & The full contract \\
Winsock & \file{ports/windows/eip_hal_windows.c}, without changes & Nothing \\
\bottomrule
\end{tabularx}
\end{table}

The platform file is short. For ESP-IDF it is approximately 190 lines with
comments, and for ClearCore approximately 130.

\section{General Rules of the Contract}

\begin{description}[font=\sffamily\bfseries\color{steel}]
  \item[Handles.] A socket handle is an \code{int}. A negative value is not
        valid, and \code{EIP_SOCK_INVALID} is $-1$. The core records which
        handles it owns, so 0 is a valid handle.
  \item[Addresses.] An IPv4 address is \code{uint8_t[4]} in network order,
        for example \code{\{192,168,1,10\}}. A port number is a plain integer
        in host order.
  \item[Timeouts.] \code{timeout_ms == 0} means ``do not wait.'' Return what
        is ready now.
  \item[Blocking waits.] Each blocking wait must keep the network stack in
        operation. On bare metal, call the stack poll inside the loop. This is
        the purpose of \code{eip_hal_yield}.
  \item[Return codes.] \code{EIP_OK} (0) or a negative \code{EIP_ERR_*},
        unless the function gives a different rule.
\end{description}

\section{Time and Scheduling}

\begin{table}[htbp]
\centering
\caption{Port functions: time, scheduling, and lock}\label{tab:halttime}
\small
\begin{tabularx}{\linewidth}{@{}P{0.34\linewidth}L@{}}
\toprule
\textbf{Function} & \textbf{Must} \\
\midrule
\code{eip_hal_millis} & Return a free-running millisecond counter that wraps at $2^{32}$. The core uses only differences. \TO{} packets are timed from it, so use a real 1~ms source if the platform has one. \\
\code{eip_hal_clock_resolution_ms} & Return the real step of the counter in milliseconds (1 for a 1~kHz source). The adapter never sets a Class~1 timeout shorter than two steps. \\
\code{eip_hal_yield} & Let other work run for a short time. On an RTOS, sleep one tick. On bare metal, run the poll of the network stack, for example \code{EthernetMgr.Refresh()}, or \code{sys_check_timeouts()} plus \code{ethernetif_input()}. \\
\code{eip_hal_lock}, \code{eip_hal_unlock} & Give mutual exclusion between the adapter context and application tasks. The lock is not nested, is held for microseconds, and is never held during a socket call. The core can read \code{eip_hal_millis} while it holds the lock. Single-threaded ports leave the two functions empty. \\
\bottomrule
\end{tabularx}
\end{table}

\begin{caution}
Use a clock with a real 1~ms step if the platform has one. On Windows,
\code{GetTickCount64} steps approximately each 16~ms. With that clock, the
adapter sent \TO{} at approximately half the rate, and the connection timed
out. \file{tests/test_hal.c} measures the step and fails a clock that is
coarser than it reports.
\end{caution}

\section{Network Interfaces}

\code{eip_hal_netif_get(index, info)} fills in the data of interface
\code{index}. Index 0 is the primary interface. These rules apply:

\begin{itemize}
  \item Return \code{EIP_ERR_ARG} after the last interface.
  \item An interface without an address still returns \code{EIP_OK}, with
        \code{ip} = 0.0.0.0. The POSIX port is an exception. It lists only
        interfaces that are up and have an IPv4 address.
  \item Set to zero all fields that you do not know.
  \item Set \code{link_known = 1} when you can report \code{link_up}. If you
        do not, the server uses ``has an address'' as ``link up, 100~Mbit/s
        full duplex.'' Fill \code{speed_mbps}, \code{full_duplex}, and
        \code{autonegotiated} with real values if you can.
\end{itemize}

The core uses the result in three ways:

\begin{enumerate}
  \item The adapter calls the function each 100~ms for \code{netif_index}. It
        fills in its Identity, TCP/IP Interface, and Ethernet Link objects.
  \item The scan reads indexes 0, 1, 2, and so on, to find broadcast addresses.
  \item The MAC address sets the default serial number.
\end{enumerate}

\section{TCP}

\begin{table}[htbp]
\centering
\caption{Port functions: TCP}\label{tab:haltcp}
\small
\begin{tabularx}{\linewidth}{@{}P{0.30\linewidth}L@{}}
\toprule
\textbf{Function} & \textbf{Must} \\
\midrule
\code{eip_hal_tcp_connect} & Connect, and wait up to \code{timeout_ms}. Return a handle or \code{EIP_SOCK_INVALID}. \\
\code{eip_hal_tcp_listen} & Listen on all addresses. Return a handle or \code{EIP_SOCK_INVALID}. \\
\code{eip_hal_tcp_accept} & \textbf{Never block.} Return a connected handle and fill \code{peer_ip}, or return \code{EIP_SOCK_INVALID} when no client waits. \\
\code{eip_hal_tcp_send} & Queue \textbf{all} \code{len} bytes. Return \code{EIP_ERR_IO} if the connection is broken, or if the bytes cannot be queued in a few hundred milliseconds. Do not block for an unlimited time. \\
\code{eip_hal_tcp_recv} & Return the number of bytes read ($>0$), 0 if nothing arrived in time, or \code{EIP_ERR_IO} after the peer closed or the connection failed. With \code{timeout_ms == 0}, return immediately. \\
\code{eip_hal_tcp_local_ip} & Return the local address of the connection. It goes into the \TO{} socket address of a Forward\_Open. \\
\code{eip_hal_tcp_close} & Close the socket and release the handle. This applies to connected and listening sockets. \\
\bottomrule
\end{tabularx}
\end{table}

For adapter sessions, set TCP\_NODELAY and keep-alive (approximately 30~s
idle, three probes). Table~\ref{tab:tcpopts} shows what the shipped ports do.

\begin{table}[htbp]
\centering
\caption{TCP options in the shipped socket layers}\label{tab:tcpopts}
\small
\begin{tabularx}{\linewidth}{@{}P{0.18\linewidth}P{0.22\linewidth}L@{}}
\toprule
\textbf{Layer} & \textbf{TCP\_NODELAY} & \textbf{Keep-alive} \\
\midrule
\file{bsd_sockets} & All connections & Accepted sockets only. The 30~s setting needs \code{TCP_KEEPIDLE}, \code{TCP_KEEPINTVL}, and \code{TCP_KEEPCNT}. \\
\file{lwip_raw} & All connections & Only in lwIP builds with \code{LWIP_TCP_KEEPALIVE}. \\
\file{netburner} & No & No. The slot of a lost peer comes back through eviction. \\
\bottomrule
\end{tabularx}
\end{table}

\section{UDP}

\begin{table}[htbp]
\centering
\caption{Port functions: UDP}\label{tab:haludp}
\small
\begin{tabularx}{\linewidth}{@{}P{0.30\linewidth}L@{}}
\toprule
\textbf{Function} & \textbf{Must} \\
\midrule
\code{eip_hal_udp_open} & Bind on all addresses. \code{local_port == 0} means any free port. \textbf{Fail if the port is already in use.} \code{EIP_UDP_BROADCAST} permits broadcast sends. \code{EIP_UDP_MULTICAST_TX} permits multicast sends with TTL 1 on the primary interface, where the stack can set them. \\
\code{eip_hal_udp_sendto} & Send one datagram. Return \code{EIP_OK} or \code{EIP_ERR_IO}. \\
\code{eip_hal_udp_recvfrom} & Receive one datagram. Return its length, 0 if nothing arrived in time, or \code{EIP_ERR_IO}. Cut a long datagram to \code{cap}. \code{ip} and \code{port} can be null. \\
\code{eip_hal_udp_close} & Close the socket and release the handle. \\
\bottomrule
\end{tabularx}
\end{table}

\begin{caution}
Do not use \code{SO_REUSEADDR}. The scanner finds a free port in UDP
2223--2231 because a bind to a port in use fails. The same failure stops an
explicit \code{local_port} of 2222 while an adapter runs on the same device.
\end{caution}

A socket bound to the any-address must receive broadcasts. The adapter
answers broadcast ListIdentity requests on UDP 44818.

\section{Wait and the Adapter Runner}

\begin{table}[htbp]
\centering
\caption{Port functions: wait and runner}\label{tab:halrun}
\small
\begin{tabularx}{\linewidth}{@{}P{0.30\linewidth}L@{}}
\toprule
\textbf{Function} & \textbf{Must} \\
\midrule
\code{eip_hal_wait} & Wait until one of the handles can be readable (or a listener has a connection), or until the timeout. Skip handles that are not valid. The function can return early. Without \code{select()}, a loop of \code{eip_hal_yield} and readiness checks is correct. \\
\code{eip_hal_adapter_start_task} & Called one time by \code{eip_adapter_start}. \textbf{RTOS:} make a task that runs \code{for (;;) eip_server_service(server, EIP_ADAPTER_IDLE_WAIT_MS);}. Give it a stack of 4--6~KB. \textbf{Bare metal:} do nothing and return \code{EIP_OK}. The application then calls \code{eip_adapter_poll}. \\
\bottomrule
\end{tabularx}
\end{table}

The server itself is static. \code{eip_server_service} needs approximately
1~KB of stack, plus the call depth of the socket layer of your platform.

\section{Check a New Port}\label{sec:checkport}

Do these checks before you connect a PLC.

\begin{procedure}{Check a new port}
\begin{enumerate}
  \item Compile with warnings on (\code{-Wall -Wextra}). The core compiles
        without warnings on the compilers of the shipped ports.
  \item If your platform can run it, run \file{tests/test_hal.c}. Make sure
        that all checks pass.
  \item On the target, start the test adapter with
        \code{eip_adapter_start(0)}.
  \item On a PC on the same network, build the POSIX version of
        \code{eiptool} (section~\ref{sec:hostbuild}).
  \item Run the commands in the listing below. Make sure that each result is
        correct.
  \item Test Class~1 from the PC. Use \code{eip_io_open} with config 150,
        output 100, input 101, and 8 bytes in each direction, as in
        \file{tests/test_loopback.c}.
  \item Let the connection run for some minutes. On the target, monitor
        \code{rx_packets}, \code{timeouts}, and the \code{rx_dropped_*}
        counters in the snapshot.
  \item On the PC, run \code{eiptool adapter}.
  \item From the target, call \code{eip_list_identity},
        \code{eip_session_open}, \code{eip_get_attribute}, and
        \code{eip_io_open} against the PC.
  \item Test with a real PLC and an EDS file that agrees with your
        assemblies.
\end{enumerate}
\end{procedure}

\begin{lstlisting}[language=shell]
./build/eiptool scan                     # finds "EIP Test Adapter"
./build/eiptool identity <target-ip>
./build/eiptool get <target-ip> 1 1 7    # Identity product name
./build/eiptool get <target-ip> 4 101 3  # produced data: 11 22 33 44 55 66 77 88
./build/eiptool get <target-ip> 0xF5 1 5 # TCP/IP interface configuration
\end{lstlisting}

If a check fails, compare your port with the POSIX port. It is the
reference, and \code{test_loopback} runs all paths through it.

\section{Rules for Changes to the Core}

If you change the core itself, obey these rules:

\begin{itemize}
  \item Keep \file{src/} free of OS calls. All OS access goes through
        \file{eip_hal.h}.
  \item Add a test with each change of behavior: \file{tests/test_core.c}
        for the codec and the adapter core, \file{tests/test_session.c} for
        scanner sessions, and \file{tests/test_hal.c} for the port contract.
  \item Use \code{/* */} comments in ASCII. Tell why, not what.
\end{itemize}
